\chapter{Aus Welt und Zeit}
\quellseite{8}

\subsection[Das entsetzliche Unglück]{Das entsetzliche Unglück des Theaterbrandes in Chicago}

durch das in wenigen Minuten Hunderte von Menschen eines grauenvollen Todes starben, ist noch in aller Munde. Die Tages- und Wochenzeitungen haben spaltenlange Berichte über die Einzelheiten gebracht, und wir brauchen nicht weiter darauf einzugehen. Das Ganze ist eine tieferschütternde Predigt. In wenig Augenblicken wurden diese Theaterbesucher aus der irdischen Lustbarkeit herausgerissen und vor den Thron des göttlichen Richters gestellt. Die Welt fragt freilich nichts nach solcher Predigt, und in kurzer Zeit wird die Aufregung sich gelegt haben und der Vorfall durch andere Ereignisse in den Hintergrund gedrängt sein. Daß doch alle Christen, insbesondere unsere Jünglinge und Jungfrauen, denen besonders die Versuchung zur Theilnahme am Weltwesen droht, diese Mahnung beherzigen und nicht vergessen möchten! Denn das heutige Theater mit seinen zum weitaus größten Theil anstößigen und unsittlichen Aufführungen ist ein Stück des sündlichen Weltwesens, dem die Christen schon in ihrer Taufe entsagt haben. Gerade auch davon gilt das Wort der Schrift: „Habt nicht lieb die Welt, noch was in der Welt ist. So jemand die Welt lieb hat, in dem ist nicht die Liebe des Vaters. Denn alles, was in der Welt ist (nämlich des Fleisches Lust und der Augen Lust und hoffärtiges Leben), ist nicht vom Vater, sondern von der Welt. Und die Welt vergehet mit ihrer Lust; wer aber den Willen Gottes thut, der bleibet in Ewigkeit“, 1 Joh. 2, 15. ff.

L. F.
